\section{Posterior experiments and reductions}
\label{fnd:section}

We represent the spin set by
\(\mathcal S=\{1,\omega,\omega^2\}\), where
\(\omega=\exp(2\pi\mathrm{i}/3)\), and write \(S_v\) for the
complex spin at a vertex \(v\). In the probabilistic recursions,
\(D\) denotes an offspring count, and
\begin{equation}\label{fnd:offspring}
 d=\E D,\qquad
 \alpha=\frac{\E[D(D-1)]}{d^2}
 =\begin{cases}1-1/d,&D\equiv d=b,\\1,&D\sim\operatorname{Pois}(d).
 \end{cases}
\end{equation}
All probabilities in the Poisson model are unconditional: the observed
tree is part of the data, and extinction is not removed from the sample
space.

\subsection{Observations, likelihoods, and posterior coordinates}

For calculations we construct the tree with ordered children.  At each
vertex, independently, we draw its offspring count and then the spins of
its children through \(K_\lambda\).  The root spin is uniform and is
independent of all offspring counts.  The resulting tree is locally
finite almost surely.  Let
\[
 O_\ell=(T_{\leq\ell},(S_v)_{|v|=\ell})
\]
denote its depth-\(\ell\) observation, initially with this auxiliary
ordering retained.  The unlabelled observation retains the rooted tree
and its boundary-spin decoration, but forgets all sibling orders.

\begin{lemma}[Removing sibling orders]\label{fnd:ordered}
The ordered and unlabelled depth observations are equivalent for inference
of the root spin.  In particular, they have the same reconstruction
advantage and the same distribution of the root posterior.
\end{lemma}

\begin{proof}
Fix a finite unlabelled boundary-decorated tree.  Conditional on any
specified root spin, the probabilities of all its distinct ordered
representatives are equal.  Indeed, a permutation of siblings permutes
independent, identically distributed descendant constructions; summing
over unobserved internal spins preserves this invariance.  The group of
such sibling permutations is transitive on the distinct representatives.
Repeated representatives caused by identical subtrees have equal
stabilizer multiplicities, so they do not alter the equality of these
probabilities.

Consequently, conditional on the unlabelled observation, the law of the
ordering is independent of the root spin.  Forgetting the ordering is a
deterministic operation, and sampling an ordering from this conditional
law is an operation that does not use the root spin.  These two operations
establish the claimed equivalence, including when the boundary is empty.
\end{proof}

An \emph{experiment} consists of a uniform spin \(S\in\mathcal S\) and
an observed datum \(Y\).  Write
\begin{equation}\label{fnd:posterior}
 p_i=\PP(S=i\mid Y),\qquad m_i=3p_i,\qquad v_i=m_i-1,
 \qquad i\in\mathcal S.
\end{equation}
Thus the likelihood vector \(m\) takes values in the compact simplex
\[
 \mathcal M=\{(m_i)_{i\in\mathcal S}:m_i\geq0,\ \textstyle\sum_i m_i=3\}.
\]
For coordinatewise expressions we use
\(\langle h\rangle=\frac13\sum_{i\in\mathcal S}h_i\).
When coordinates are written as numerical triples indexed by $0,1,2$,
coordinate $i\in\{0,1,2\}$ corresponds to the spin label $\omega^i$.
The complex posterior and its two basic invariants are
\begin{equation}\label{fnd:complex}
 \zeta=\E[S\mid Y]=\langle i m_i\rangle,\qquad
 x=|\zeta|^2,\qquad y=\operatorname{Re}(\zeta^3).
\end{equation}
Here the \(i\) inside the coordinate average is the spin label, not the
imaginary unit.  Fourier inversion on the three-element group gives
\begin{equation}\label{fnd:fourier}
 v_i=\zeta\overline{i}+\overline{\zeta}i
     =2\operatorname{Re}(\zeta\overline{i}),\qquad
 \langle v^2\rangle=2x,\qquad \langle v^3\rangle=2y.
\end{equation}
More precisely, each of the three coordinates satisfies
\begin{equation}\label{fnd:cubic}
 v_i^3=3xv_i+2y,\qquad
 \prod_{i\in\mathcal S}(1+t v_i)=1-3t^2x+2t^3y.
\end{equation}
To verify the first identity, expand
\((\zeta\overline{i}+\overline{\zeta}i)^3\) and use \(i^3=1\).
The second follows from \(\sum_i v_i=0\),
\(\sum_i v_i^2=6x\), and \(\prod_i v_i=2y\).
The convex-hull constraint on \(\zeta\), together with
\(\prod_i m_i\geq0\), gives
\begin{equation}\label{fnd:invariant-bounds}
 0\leq x\leq1,\qquad |y|\leq x^{3/2},\qquad 1-3x+2y\geq0.
\end{equation}
For later use these also imply, when \(x>0\),
\begin{equation}\label{fnd:ratio-domain}
 -\frac12\leq r:=\frac yx\leq1,\qquad
 r^2\leq x\leq\frac1{3-2r}.
\end{equation}
Indeed \(|r|\leq\sqrt{x}\leq1\), and the last inequality in
Equation~\eqref{fnd:invariant-bounds} gives the upper bound on \(x\).
Combining the two bounds on \(x\) gives
\(0\leq1-3r^2+2r^3=(1-r)^2(1+2r)\), and hence \(r\geq-1/2\).

The marginal law of a posterior likelihood vector is always balanced:
\(\E m_i=1\).  Bayes' formula states that its conditional law under
input \(i\) is tilted by \(m_i\).  More generally, for a bounded
measurable function \(g\) of the original datum,
\begin{equation}\label{fnd:bayes-tilt}
 \E[g(Y)\mid S=i]=\E[m_i(Y)g(Y)].
\end{equation}
Conversely, every probability law \(Q\) on \(\mathcal M\) satisfying
\(\int m_i\,Q(dm)=1\) defines an experiment by taking
\(Q(dm\mid S=i)=m_iQ(dm)\).  Its posterior likelihood vector is
precisely the observed \(m\).  Thus balanced laws on \(\mathcal M\)
are sufficient to describe the experiments used below.

The posterior vector also loses no inferential information from an
existing experiment.  In fact
\(\PP(S=i\mid Y,m)=m_i/3=\PP(S=i\mid m)\), so \(S\) and \(Y\)
are conditionally independent given \(m\).  The original datum can
therefore be conditionally simulated from \(m\) without knowing \(S\).
All observation spaces here are standard Borel, so the required
conditional kernels exist.

Our experiments are invariant under every permutation of the three spin
labels.  In complex coordinates these permutations are generated by
\(\zeta\mapsto\omega\zeta\) and
\(\zeta\mapsto\overline{\zeta}\).  This invariance is retained by
all the tree recursions and limits below.

\begin{lemma}[Total variation and the posterior second moment]
\label{fnd:norms}
For every posterior vector,
\begin{equation}\label{fnd:tv-pointwise}
 \frac{|\zeta|}{\sqrt3}
 \leq\left\|p-\operatorname{Unif}(\mathcal S)\right\|_{\mathrm{TV}}
 \leq\frac23|\zeta|.
\end{equation}
Consequently, writing \(\mu_\ell=\E x_\ell\) for the depth experiment
and \(a_\ell\) for its mean total variation advantage,
\begin{equation}\label{fnd:tv-moment}
 \frac{\mu_\ell}{\sqrt3}\leq a_\ell
 \leq\frac23\sqrt{\mu_\ell}.
\end{equation}
In particular, vanishing advantage is equivalent to vanishing posterior
second moment.
\end{lemma}

\begin{proof}
For any three real numbers with zero sum, half the sum of their absolute
values equals their largest absolute value.  Applying this to \(v\)
gives \(\|p-\operatorname{Unif}(\mathcal S)\|_{\mathrm{TV}}
=\max_i|v_i|/3\).  If \(M=\max_i|v_i|\), their squared sum is at
most \(2M^2\): after choosing the coordinate of absolute value \(M\),
the other two have the opposite sign and their absolute values sum to
\(M\).  Equation~\eqref{fnd:fourier} therefore yields
\(6x\leq2M^2\).  The same equation gives
\(M\leq2|\zeta|\), proving Equation~\eqref{fnd:tv-pointwise}.
Now use \(|\zeta|^2\leq|\zeta|\) and Cauchy--Schwarz to obtain
Equation~\eqref{fnd:tv-moment}.
\end{proof}

\subsection{The exact posterior recursion}

Suppose \(Y\) is a child experiment with likelihood vector \(m\).
The child spin has a uniform marginal, and the conditional law of \(Y\)
given parent spin \(i\), relative to the child marginal law, has density
\begin{equation}\label{fnd:edge}
 m_{e,i}=\sum_{k\in\mathcal S}K_\lambda(i,k)m_k
        =1+\lambda(m_i-1)=1+\lambda v_i.
\end{equation}
Thus the datum retains its original marginal law under the uniform
parent prior.  The edge operation sends
\begin{equation}\label{fnd:edge-moments}
 \zeta\longmapsto\lambda\zeta,\qquad
 x\longmapsto\lambda^2x,\qquad y\longmapsto\lambda^3y.
\end{equation}

To combine \(n\) experiments that are independent conditional on a
common input, let \(Q_1,\ldots,Q_n\) be their marginal likelihood laws.
Under the \emph{product marginal law} \(Q_1\otimes\cdots\otimes Q_n\),
set
\begin{equation}\label{fnd:product}
 W_i=\prod_{j=1}^n m_{i,j},\qquad
 Z=\langle W\rangle=\frac13\sum_{i\in\mathcal S}W_i,\qquad
 m_i^{\mathrm{out}}=\frac{W_i}{Z}\quad(Z>0).
\end{equation}
Conditional on input \(i\), the joint datum has density \(W_i\)
relative to that product law.  Averaging over the input gives the
marginal density \(Z\), and Bayes' formula gives the last expression
in Equation~\eqref{fnd:product}.  In particular,
\(\int Z\,d(Q_1\otimes\cdots\otimes Q_n)=1\).
The value of \(m^{\mathrm{out}}\) on \(\{Z=0\}\) may be set equal
to \((1,1,1)\); this set has zero probability after the tilt.
When \(n=0\), empty products are one, \(Z=1\), and the output is
the uninformative vector \((1,1,1)\).

\begin{proposition}[Density evolution]\label{fnd:recursion}
Let \(Q_\ell\) be the marginal likelihood law of the depth-\(\ell\)
root experiment.  Then
\begin{equation}\label{fnd:law-recursion}
 Q_{\ell+1}=\mathcal R_{\lambda,D}(Q_\ell),\qquad
 Q_0=\frac13\sum_{i\in\mathcal S}\delta_{3\mathbf e_i},
\end{equation}
where \(\mathbf e_i\) is the coordinate unit vector and, for every
bounded measurable \(f\),
\begin{equation}\label{fnd:operator}
 \int f\,d\mathcal R_{\lambda,D}(Q)
 =\sum_{n\geq0}\PP(D=n)
   \int Z f(m^{\mathrm{out}})\,dQ^{\otimes n}.
\end{equation}
In this formula Equation~\eqref{fnd:product} is applied to
\(m_{i,j}=1+\lambda(\widetilde m_{i,j}-1)\), where
\(\widetilde m_{\cdot,j}\) are the independent vectors sampled from
\(Q\).  Each conditional integral has total mass one.  The operator
preserves balanced laws and permutation symmetry.
\end{proposition}

\begin{proof}
Conditional on the offspring count and the parent spin, the child
subtree observations are independent copies of the depth-\(\ell\)
experiment passed through one edge.  Equations~\eqref{fnd:edge} and
\eqref{fnd:product} give their combination.  The count is observed and
independent of the parent spin, so it is mixed with its original law
in Equation~\eqref{fnd:operator}.  For each \(i\), the expectation of
the output coordinate is \(\int W_i\,dQ^{\otimes n}=1\), proving
balance.  Relabelling the input coordinates commutes with every
operation.  At depth zero the root itself is observed, giving \(Q_0\).
Random subtree shapes are already included in each child datum; no
additional conditioning or averaging convention is required.
\end{proof}

We say that one experiment is \emph{degraded} from another when its
datum is obtained from the latter by a random kernel independent of
the input spin.  If \(Y'\) is degraded from \(Y\), their posterior
likelihoods satisfy
\begin{equation}\label{fnd:conditional-jensen}
 m(Y')=\E[m(Y)\mid Y'],\qquad
 \E\Phi(m(Y'))\leq\E\Phi(m(Y))
\end{equation}
for every integrable convex function \(\Phi\) of the posterior vector.
This is the tower property followed by conditional Jensen.  In
particular it applies to total variation, \(x\), and \(x^2\): total
variation is the convex function \(\max_i|m_i-1|/3\), while
\(\zeta\) is linear in \(m\) and its squared and fourth powers of
the modulus are convex.  Applying degradation
kernels separately to conditionally independent branches shows that
edge transmission and branch composition preserve this order.

\subsection{Convergence, fixed points, and extinction}

\begin{proposition}[The limiting posterior]\label{fnd:limit}
On a common realization of the tree and spins, the root posterior
likelihood vectors \(m^{(\ell)}\) converge almost surely to a vector
\(m^{(\infty)}\).  Their laws converge weakly to a symmetric balanced
law \(Q_\infty\) satisfying
\begin{equation}\label{fnd:fixed-point}
 Q_\infty=\mathcal R_{\lambda,D}(Q_\infty).
\end{equation}
The advantages \(a_\ell\) and moments \(\mu_\ell=\E x_\ell\) are
nonincreasing and converge to their corresponding expectations under
\(Q_\infty\).  This holds throughout \(-1/2\leq\lambda\leq1\),
including parameters for which edge probabilities vanish.
\end{proposition}

\begin{proof}
The observations \(O_\ell\) themselves do not generate a decreasing
sequence of sigma-fields.  Instead put
\[
 \mathcal F_\ell=\sigma(O_k:k\geq\ell).
\]
Given \(O_\ell\), the shapes and spins below its boundary are generated
independently of the root spin.  This statement includes every finite
collection of subsequent depth observations, and hence their generated
sigma-field.  Therefore
\[
 \E[\mathbf1_{\{S_\rho=i\}}\mid\mathcal F_\ell]
 =\PP(S_\rho=i\mid O_\ell)=\tfrac13m_i^{(\ell)}.
\]
The reverse martingale convergence theorem, applied to these bounded
indicators and the decreasing fields \(\mathcal F_\ell\), gives
almost-sure and \(L^1\) convergence.  The limit is the root posterior
given \(\bigcap_\ell\mathcal F_\ell\).  Every continuous function of
the posterior converges in expectation by bounded convergence on the
compact simplex.  Moreover
\(m^{(\ell+1)}=\E[m^{(\ell)}\mid\mathcal F_{\ell+1}]\), so
conditional Jensen proves the stated monotonicity.  Lemma~\ref{fnd:ordered}
transfers these conclusions to the required unlabelled observations.

It remains to pass Equation~\eqref{fnd:law-recursion} to the limit.
For a fixed count \(n\) and a bounded continuous \(f\), the integrand
\(Zf(m^{\mathrm{out}})\) in Equation~\eqref{fnd:operator} is continuous
where \(Z>0\).  Give it the value zero where \(Z=0\).  This extension
is continuous, since
\[
 |Zf(m^{\mathrm{out}})|\leq\|f\|_\infty Z
\]
and \(Z\) is continuous and nonnegative.  Weak convergence of the
finite product laws therefore gives convergence of the integral for
each fixed \(n\), even at the coloring endpoint.  The tail of the
count mixture is uniformly bounded by
\begin{equation}\label{fnd:count-tail}
 \left|\sum_{n>M}\PP(D=n)
     \int Zf(m^{\mathrm{out}})\,dQ^{\otimes n}\right|
 \leq\|f\|_\infty\PP(D>M),
\end{equation}
because each conditional tilted law has mass one.  First truncate the
count, then pass to the weak limit, and finally let \(M\to\infty\).
Thus \(\mathcal R_{\lambda,D}\) is weakly continuous on balanced
laws.  Since both \(Q_\ell\) and \(Q_{\ell+1}\) converge to
\(Q_\infty\), Equation~\eqref{fnd:fixed-point} follows.
\end{proof}

Unless a different experiment is specified, we henceforth use
\begin{equation}\label{fnd:limit-moments}
 \mu=\E x,\qquad \nu=\E x^2,\qquad \eta=\E y
\end{equation}
for the limiting experiment.  Its law is trivial precisely when
\(\mu=0\), since \(x=0\) is equivalent to \(m=(1,1,1)\).
By Lemma~\ref{fnd:norms} and Proposition~\ref{fnd:limit},
non-reconstruction is equivalent to \(\mu=0\).

\begin{lemma}[Logarithms at interior channel parameters]
\label{fnd:log-integrability}
If \(-1/2<\lambda<1\), every fixed point of
\(\mathcal R_{\lambda,D}\) has strictly positive coordinates almost
surely.  Every polynomial in its coordinates and their logarithms is
integrable under either offspring law in Equation~\eqref{fnd:offspring}.
\end{lemma}

\begin{proof}
For every \(m\in\mathcal M\), its edge coordinates lie in
\([c_-,c_+]\), where
\[
 c_- =\min(1-\lambda,1+2\lambda)>0,\qquad
 c_+ =\max(1-\lambda,1+2\lambda).
\]
At count \(n\), both \(W_i\) and \(Z\) lie in
\([c_-^n,c_+^n]\).  Hence
\begin{equation}\label{fnd:log-bound}
 |\log m_i^{\mathrm{out}}|\leq n\log(c_+/c_-).
\end{equation}
The count has its original law in the tilted output construction,
and either specified law has finite moments of every order.  The
coordinates themselves lie in \([0,3]\).  The bound proves the
assertion, and also bounds each centered log coordinate
\(\log m_i-\langle\log m\rangle\) by
\(2n\log(c_+/c_-)\).  At count zero all these logarithms are zero.
No claim of this kind is needed at \(\lambda=-1/2\).
\end{proof}

\begin{lemma}[The extinction factor]\label{fnd:survival}
Let \(s=1\) in the regular model and let \(s\) be the survival
probability in the Poisson model.  For the limiting experiment,
\begin{equation}\label{fnd:survival-moments}
 \mu^2\leq s\nu,\qquad |\eta|\leq\sqrt{\mu\nu}.
\end{equation}
In the Poisson model with \(d>1\), \(0<s<1\) and
\begin{equation}\label{fnd:survival-equation}
 -\log(1-s)=ds,\qquad
 s\leq2(1-1/d).
\end{equation}
In particular, at \(d\lambda^2=1\),
\(s\leq2(1-\lambda^2)\).
\end{lemma}

\begin{proof}
On extinction the boundary is empty at all sufficiently large depths.
An observation with empty boundary consists only of a tree shape, which
is independent of the root spin; its posterior is uniform.  Thus
\(x_\infty=0\) on extinction.  Cauchy--Schwarz gives
\(\mu=\E[x_\infty\mathbf1_{\{\mathrm{survival}\}}]
\leq\sqrt{s\nu}\).  Also
\(|\eta|\leq\E x^{3/2}\leq\sqrt{\mu\nu}\) by
Equation~\eqref{fnd:invariant-bounds}.

For completeness, the extinction probability is the increasing limit
of the probabilities of an empty generation.  The offspring generating
function gives it as the smallest solution \(q\in[0,1]\) of
\(q=\exp(d(q-1))\).  Strict convexity and the derivative \(d>1\)
at \(q=1\) show that this smallest solution belongs to \((0,1)\).
Setting \(q=1-s\) proves the first identity in
Equation~\eqref{fnd:survival-equation}.  For \(0<s<1\),
\[
 -\log(1-s)=\sum_{k\geq1}\frac{s^k}{k}
 \geq\sum_{k\geq1}\frac{s^k}{2^{k-1}}
 =\frac{s}{1-s/2},
\]
where \(2^{k-1}\geq k\).  Dividing by \(s\) and using the
survival equation gives \(d\geq(1-s/2)^{-1}\), equivalently the
second inequality in Equation~\eqref{fnd:survival-equation}.
\end{proof}

\subsection{Reconstruction above the second-moment threshold}

\begin{proposition}[A linear estimator]\label{fnd:supercritical}
In either model, \(d\lambda^2>1\) implies reconstruction, for either
sign of \(\lambda\).
\end{proposition}

\begin{proof}
The hypothesis implies \(\lambda\neq0\).  Define the observable,
order-independent estimator
\begin{equation}\label{fnd:linear-estimator}
 L_\ell=(d\lambda)^{-\ell}\sum_{|u|=\ell}S_u,
 \qquad L_0=S_\rho.
\end{equation}
The sum is zero on an empty boundary.  Since
\(\E[S_{\mathrm{child}}\mid S_{\mathrm{parent}}]
=\lambda S_{\mathrm{parent}}\), induction gives
\(\E[L_\ell\mid S_\rho]=S_\rho\).

Let \(L_\ell^{(j)}\) denote the corresponding estimator in the
subtree rooted at child \(j\).  Then
\(L_{\ell+1}=(d\lambda)^{-1}\sum_{j=1}^D L_\ell^{(j)}\).
Conditional on \(S_\rho\) and \(D\), these variables are independent,
and each has mean \(\lambda S_\rho\).  The diagonal terms in the
squared modulus sum to \(d\E|L_\ell|^2\); each ordered off-diagonal
pair has expectation \(\lambda^2\).  Consequently, with
\(q=(d\lambda^2)^{-1}<1\),
\begin{equation}\label{fnd:estimator-recurrence}
 \E|L_{\ell+1}|^2
 =q\E|L_\ell|^2+\alpha,\qquad
 \E|L_\ell|^2=q^\ell+\alpha\frac{1-q^\ell}{1-q}
 \leq1+\frac{\alpha}{1-q}=:C.
\end{equation}
The factor \(\alpha\) is the normalized factorial moment, rather
than the normalized second moment, because it counts distinct branches.

The correlation with the root equals one.  Conditioning on the observed
experiment and applying Cauchy--Schwarz gives
\[
 1=\E[L_\ell\overline{S_\rho}]
   =\E[L_\ell\overline{\zeta_\ell}],\qquad
 1\leq\E|L_\ell|^2\,\E|\zeta_\ell|^2.
\]
Thus \(\mu_\ell\geq C^{-1}\) at every depth, and
Equation~\eqref{fnd:tv-moment} gives
\(a_\ell\geq(\sqrt3 C)^{-1}>0\).  The limiting advantage exists by
Proposition~\ref{fnd:limit}, so it is positive.  In particular, this
argument uses the original Poisson experiment, without conditioning
on survival.
\end{proof}

\subsection{Decreasing the channel strength or the offspring mean}

\begin{proposition}[Channel and offspring degradation]
\label{fnd:degradation}
At every finite depth, each of the following changes produces a
degraded root experiment:
\begin{enumerate}
 \item replacing \(\lambda\) by \(c\lambda\), where \(0\leq c\leq1\);
 \item replacing a regular offspring count \(b_+\) by an integer
       \(2\leq b_-\leq b_+\);
 \item replacing a Poisson offspring mean \(d_+\) by
       \(0<d_-\leq d_+\).
\end{enumerate}
Consequently each change can only decrease the expected total variation
advantage and the moments \(\E x\) and \(\E x^2\), both at finite
depth and in their limits.
\end{proposition}

\begin{proof}
The channel admits a group-increment construction.  On each edge \(e\)
take an independent \(G_e\in\mathcal S\) with
\[
 \PP(G_e=1)=\frac{1+2\lambda}{3},\qquad
 \PP(G_e=\omega)=\PP(G_e=\omega^2)=\frac{1-\lambda}{3},
\]
and set \(S_v=S_\rho\prod_{e\in[\rho,v]}G_e\).
The increment law is invariant under conjugation and has mean
\(\E G_e=\lambda\).  Independently generate increments \(H_e\)
with the same form of law at parameter \(c\).  A conjugation-invariant
probability law on \(\mathcal S\) is uniquely determined by its real
mean: if its mass at 1 is \(p\), that mean is \((3p-1)/2\).
The independent products \(G_eH_e\) are conjugation-invariant and
have mean \(c\lambda\), so they have exactly the channel increment
law for \(c\lambda\).

Given only the observed tree through depth \(\ell\) and its boundary
spins, generate the \(H_e\)'s on its edges and replace each boundary
spin by
\[
 S'_v=S_v\prod_{e\in[\rho,v]}H_e.
\]
This kernel requires neither the root spin nor any unobserved internal
spins, and the output has exactly the depth experiment for parameter
\(c\lambda\).  Shared path increments give the required joint law
of the transformed boundary, not merely its individual marginals.

For regular trees, in the ordered representation retain the first
\(b_-\) children at each retained vertex and discard the rest of
their subtrees.  The retained boundary and tree have the smaller regular
experiment.  For Poisson trees, independently keep each child of a
retained vertex with probability \(p=d_-/d_+\).  If \(D_+\) is its
original count and \(D_-\) its retained count, then
\[
 \E z^{D_-}=\E(1-p+pz)^{D_+}
             =\exp(d_-(z-1)).
\]
Independence across vertices and independence from the edge increments
show that the retained observation has the Poisson broadcast law with
mean \(d_-\).  These operations are performed only on the known tree
and observed boundary; neither requires a survival event.

By Lemma~\ref{fnd:ordered}, an unlabelled observation can first be
randomly ordered without accessing the root, and the final ordering
can then be forgotten.  All the preceding kernels are therefore valid
for the unlabelled model as well.  Their comparison of the stated
quantities follows from Equation~\eqref{fnd:conditional-jensen}, and
passes to the limits by Proposition~\ref{fnd:limit}.
\end{proof}

At \(\lambda=0\), every positive-depth boundary is independent of
the root conditional on the tree, so the advantage is zero.  For each
fixed \(d\), Proposition~\ref{fnd:degradation} reduces the remaining
subcritical channel parameters of either sign to their admissible
equality parameters.  Negative equality is admissible exactly when
\(d\geq4\); for smaller branching parameters, non-reconstruction at
the coloring channel \(\lambda=-1/2\) and branching parameter four
will imply all the needed negative-channel statements by the same
proposition.
